\chapter{群}
    \section{剰余類}
        \subsection{定義}
            $G$を群とし，$H$をその任意の部分群とする。
            $g\in G$を任意にとる。
            \par
            「$G$に於ける$H$の左剰余類 (the left coset of $H$ in $G$)」とは次の集合族のことを言う。
            \[ \setComprehension*{gH}{g \in G} = \setComprehension*{\setComprehension*{gh}{h \in H}}{g \in G} \]
            また，「$G$に於ける$H$の右剰余類」とは次の集合族のことを言う。
            \[ \setComprehension*{Hg}{g \in G} = \setComprehension*{\setComprehension*{hg}{h \in H}}{g \in G} \]
            特に，$gH\;(g \in G)$を「$G$に於ける$g$に関する$H$の左剰余類 (the left coset of $H$ in $G$ with respect to $g$)」と呼び，$Hg\;(g \in G)$を「$G$に於ける$g$に関する$H$の右剰余類」と呼ぶ。
        \subsection{群\texorpdfstring{$G$}{G}の部分群\texorpdfstring{$H$}{H}の全ての剰余類は互いに素である}
            \begin{proof}
                \quad\par
                左剰余類について示す。
                右剰余類についても同様にして示せる。
                \par
                命題の対偶「2つの左剰余類$R=gH,R'=g'H\;(g\in G)$の共通部分が非空ならば$R=R'$」を示す。
                $R$と$R'$の共通部分を$gh = g'h'\;(h,h'\in H)$とすると$g' = ghh'^{-1}$だから$\forall h'' \in H,(R' \ni) g'h'' = ghh'^{-1}h'' \in gH = R\;(\because\;hh'^{-1}h'' \in H)$であり$R'\subseteq R$となる。
                同様に$R' \supseteq R$も示せるから$R=R'$となる。
            \end{proof}
        \subsection{2つの元が同じ剰余類に属する必要十分条件}
            $G$を群とし，$H$をその部分群とする。
            $s_1,s_2\in G$が$H$の同じ左剰余類に属する必要十分条件は$s_1^{-1}s_2\in H$である。
            また，同じ右剰余類に属する必要十分条件は$s_1s_2^{-1}\in H$である。
            \begin{proof}
                \quad\par
                左剰余類について示す。
                右剰余類についても同様にして示せる。\\
                ($\Rightarrow$)
                \par
                仮定より$\exists g \in G,h_1,h_2 \in H \st s_1 = gh_1,s_2 = gh_2$。
                よって$s_1^{-1}s_2 = h_1^{-1}g^{-1}gh_2 = h_1^{-1}h_2 \in H$\\
                ($\Leftarrow$)
                \par
                仮定より$\exists h \in H \st s_1^{-1}s_2 = h$。
                直前に示した定理より$H$の剰余類は$G$を分割するから，$\exists g\in G,h_1\in H \st s_1 = gh_1 (\in gH)$。
                よって$s_2 = s_1h = gh_1h \in gH\;(\because\; h_1h \in H)$
            \end{proof}
        \section{\texorpdfstring{$\mathbb{Z}_p$}{Zp} (\texorpdfstring{$p: $}{p: }素数)に乗法群としての位数\texorpdfstring{$r$}{r}の元があれば\texorpdfstring{$p-1$}{p-1}は\texorpdfstring{$r$}{r}の倍数}
            \begin{proof}
                \quad\par
                論題の元を$a$とする。
                $r$は$a^r \equiv_p 1$となる最小の自然数であるから，$k\in\naturalNumbers$に対して$a^k \equiv_p 1$となるのは$k$が$r$の倍数であるとき，かつそのときに限る。
                Fermatの小定理から$a^{p-1} \equiv_p 1$であるので，$p-1$は$r$の倍数である。
            \end{proof}